% CONTINUES: rational points on the circle $x^2+y^2=1$, via lines of rational slope through $(-1,0)$.
% --- page 62 ---

The line through $(-1,0)$ of slope $t$ is $y=t(x+1)$. Substituting into the circle gives
\[
    \bigl(t(x+1)\bigr)^2+x^2=1,
\]
which expands to the quadratic
\[
    x^2(1+t^2)+2t^2 x+(t^2-1)=0.
\]
One root is the known intersection $x=-1$. Factoring,
\[
    (x+1)\bigl((1+t^2)x-(1-t^2)\bigr)
    =x^2(1+t^2)+2t^2 x+(t^2-1),
\]
so the other root is
\[
    x=\frac{1-t^2}{1+t^2}.
\]
The corresponding $y$-coordinate is
\[
    y=t(x+1)=t\left(\frac{1-t^2}{1+t^2}+1\right)
    =\frac{2t}{1+t^2}.
\]
Thus every rational slope produces a rational point, and every rational point other than $(-1,0)$ arises this way:
\[
    \left\{
        \left(\frac{1-t^2}{1+t^2},\,\frac{2t}{1+t^2}\right)
        : t\in\mathbb{Q}
    \right\}
    \cup\{(-1,0)\}.
\]
Sending $t\to\infty$ (the vertical line through $(-1,0)$) returns the missing point: the two coordinates tend to $-1$ and $0$.

\subsection{The group law}

The name is historical. Arc lengths of ellipses lead to elliptic integrals, and elliptic integrals lead to the curves below.

\begin{keybox}
\begin{definition}[Elliptic curve]
For us, an elliptic curve is an equation
\[
    y^2=x^3+ax+b,
\]
where $a,b\in\R$ satisfy $4a^3+27b^2\neq 0$. The condition ensures that the cubic has no repeated root, so the curve is nonsingular.
\end{definition}
\end{keybox}

% --- page 63 ---

Over $\R$ the real locus has one of the two shapes in Figure~\ref{fig:ec-real}, according to the number of real roots of the cubic. A cross marks a root.

\begin{figure}[H]
\centering
\begin{tikzpicture}[x=0.78cm,y=0.78cm,line cap=round,line join=round]
    % one real component: y^2 = x^3 - x + 1, root about -1.325
    \begin{scope}
        \draw[->] (-2.15,0) -- (2.15,0);
        \draw (0,-1.85) -- (0,1.85);
        \draw[domain=-1.32472:1.85,samples=240,smooth] plot (\x,{sqrt(max(0,\x*\x*\x-\x+1))});
        \draw[domain=-1.32472:1.85,samples=240,smooth] plot (\x,{-sqrt(max(0,\x*\x*\x-\x+1))});
        \node[fill=white,inner sep=1pt] at (-1.325,0) {$\times$};
    \end{scope}
    % two components: y^2 = x^3 - 4x, roots -2,0,2
    \begin{scope}[xshift=6.6cm]
        \draw[->] (-2.7,0) -- (3.55,0);
        \draw (1,-1.85) -- (1,1.85);
        \draw[domain=-2:0,samples=240,smooth] plot (\x,{sqrt(max(0,\x*\x*\x-4*\x))});
        \draw[domain=-2:0,samples=240,smooth] plot (\x,{-sqrt(max(0,\x*\x*\x-4*\x))});
        \draw[domain=2:2.85,samples=240,smooth] plot (\x,{sqrt(max(0,\x*\x*\x-4*\x))});
        \draw[domain=2:2.85,samples=240,smooth] plot (\x,{-sqrt(max(0,\x*\x*\x-4*\x))});
        \node[fill=white,inner sep=1pt] at (-2,0) {$\times$};
        \node[fill=white,inner sep=1pt] at (0,0) {$\times$};
        \node[fill=white,inner sep=1pt] at (2,0) {$\times$};
    \end{scope}
\end{tikzpicture}
\caption{Real locus of a short Weierstrass equation. Left: $y^2=x^3-x+1$, one real root. Right: $y^2=x^3-4x$, three real roots, drawn as an oval and an unbounded component, with a vertical line in the gap between them.}
\label{fig:ec-real}
\end{figure}

This shape of equation is called the \emph{Weierstrass form}. It is available whenever $\operatorname{char} K\neq 2,3$. In the remaining characteristics one uses
\begin{align*}
    \operatorname{char} K=2:
    &\quad
    \begin{cases}
        y^2+cy=x^3+ax+b,\\
        y^2+xy=x^3+ax^2+b,
    \end{cases}\\[0.4em]
    \operatorname{char} K=3:
    &\quad
    y^2=x^3+ax^2+bx+c.
\end{align*}
The general Weierstrass equation, over any field, is
\[
    y^2+a_1 xy+a_3 y=x^3+a_2 x^2+a_4 x+a_6.
\]

An elliptic curve comes with a point at infinity. That extra point is what makes it possible to put an abelian group structure on the curve.

Adding two points by adding coordinates does not work. If $(x_1,y_1)$ and $(x_2,y_2)$ lie on $E$, the pair $(x_1+x_2,\,y_1+y_2)$ need not.

% --- page 64 ---

A second attempt uses the geometry of a line. Let $P$ and $Q$ be points of $E$, and let the line through them meet $E$ again at $R$, as in Figure~\ref{fig:ec-naive}.

\begin{figure}[H]
\centering
\begin{tikzpicture}[x=0.72cm,y=0.72cm,line cap=round,line join=round]
    \draw[->] (-2.7,0) -- (4.15,0);
    % y-axis in the gap (0,2); x=0 is a root, with a vertical tangent
    \draw[gray!55] (0,-4.3) -- (0,4.3);
    \draw[domain=-2:0,samples=240,smooth] plot (\x,{sqrt(max(0,\x*\x*\x-4*\x))});
    \draw[domain=-2:0,samples=240,smooth] plot (\x,{-sqrt(max(0,\x*\x*\x-4*\x))});
    \draw[domain=2:3.30,samples=240,smooth] plot (\x,{sqrt(max(0,\x*\x*\x-4*\x))});
    \draw[domain=2:3.30,samples=240,smooth] plot (\x,{-sqrt(max(0,\x*\x*\x-4*\x))});
    % P, Q, R lie on y^2=x^3-4x; chord y+0.397=0.881(x+1.98) ends on R
    \coordinate (P) at (-1.98,-0.397);
    \coordinate (Q) at (-0.30,1.083);
    \coordinate (R) at (3.056,4.040);
    \draw (-2.32,-0.696) -- (R);
    \draw[->] (-0.72,0.713) -- (-0.48,0.924);
    \fill (P) circle (1.6pt);
    \node[anchor=north east] at (-2.08,-0.58) {$P$};
    \fill (Q) circle (1.6pt);
    \node[above=3pt] at (Q) {$Q$};
    \fill (R) circle (1.6pt);
    \node[above left=1pt] at (R) {$R$};
\end{tikzpicture}
\caption{The chord through $P$ and $Q$ on $y^2=x^3-4x$ meets the curve again at $R$. Declaring $P+Q=R$ does not give a group law.}
\label{fig:ec-naive}
\end{figure}

\begin{remark}
Taking $P+Q$ to be this third point $R$ fails. The three intersection points are symmetric, so the same rule would also require $Q+R=P$ and $P+R=Q$, and those identities do not hold. The proposal collapses further to $2P=2Q=2R=\mathcal{O}_E$, i.e.\ every point would have order dividing $2$. The notes mark the attempt with $(\times)$; it also fails.
\end{remark}

The law that does work reflects the third intersection across the $x$-axis. Draw the line through $P$ and $Q$, let $R$ be its third meeting with $E$, and let $R'$ be the reflection of $R$ in the $x$-axis (Figure~\ref{fig:ec-sum}). Define
\[
    P+Q=R'.
\]

\begin{figure}[H]
\centering
\begin{tikzpicture}[x=0.78cm,y=0.78cm,line cap=round,line join=round]
    \draw[->] (-2.55,0) -- (4.05,0);
    % y-axis in the gap (0,2); x=0 is a root, with a vertical tangent
    \draw[gray!55] (0,-3.25) -- (0,3.15);
    \draw[domain=-2:0,samples=240,smooth] plot (\x,{sqrt(max(0,\x*\x*\x-4*\x))});
    \draw[domain=-2:0,samples=240,smooth] plot (\x,{-sqrt(max(0,\x*\x*\x-4*\x))});
    \draw[domain=2:3.25,samples=240,smooth] plot (\x,{sqrt(max(0,\x*\x*\x-4*\x))});
    \draw[domain=2:3.25,samples=240,smooth] plot (\x,{-sqrt(max(0,\x*\x*\x-4*\x))});
    % P, Q, R on y^2=x^3-4x; R' is the reflection of R
    \coordinate (P) at (-1.800,-1.170);
    \coordinate (Q) at (-0.040,0.400);
    \coordinate (R) at (2.635,2.786);
    \coordinate (Rp) at (2.635,-2.786);
    \draw (-2.15,-1.482) -- (R);
    \draw[dashed] (R) -- (Rp);
    \fill (P) circle (1.6pt);
    \node[anchor=north] at (-1.80,-1.32) {$P$};
    \fill (Q) circle (1.6pt);
    \node[anchor=east] at (-0.28,0.72) {$Q$};
    \fill (R) circle (1.6pt);
    \node[anchor=south west] at (2.78,2.90) {$R$};
    \fill (Rp) circle (1.6pt);
    \node[anchor=west] at (2.95,-2.55) {$R'=P+Q$};
\end{tikzpicture}
\caption{Chord through $P$ and $Q$ on $y^2=x^3-4x$, meeting the curve again at $R$. The dashed segment reflects $R$ across the $x$-axis to $R'$, and the sum is $P+Q=R'$.}
\label{fig:ec-sum}
\end{figure}

\par
\ifdim\dimexpr\pagegoal-\pagetotal<5\baselineskip
  \newpage
\fi
\begin{example}[Adding two distinct points]
\label{ex:pq-sum}
Take $y^2=x^3-15x+18$, with $P=(7,16)$ and $Q=(1,2)$. Both points lie on the curve:
\[
    16^2=256=7^3-15\cdot 7+18,
    \qquad
    2^2=4=1^3-15\cdot 1+18.
\]
The slope of the chord is
\[
    \frac{16-2}{7-1}=\frac{14}{6}=\frac{7}{3},
\]
so the line is
\[
    y-2=\frac{7}{3}(x-1),
    \qquad
    y=\frac{7}{3}x-\frac{1}{3}.
\]
Substitute into the curve:
\[
    \left(\frac{7}{3}x-\frac{1}{3}\right)^2=x^3-15x+18.
\]
Clearing the denominator $9$,
\begin{align*}
    (7x-1)^2&=9x^3-135x+162,\\
    49x^2-14x+1&=9x^3-135x+162,\\
    0&=9x^3-49x^2-121x+161.
\end{align*}
Dividing by $9$ gives the monic cubic whose roots are the $x$-coordinates of the three intersections:
\begin{equation}
    x^3-\frac{49}{9}x^2-\frac{121}{9}x+\frac{161}{9}=0.
    \label{eq:pq-cubic}
\end{equation}
% --- page 65 ---
The roots include $x=7$ and $x=1$. If $\lambda$ is the third, then
\[
    \lambda+1+7=\frac{49}{9},
    \qquad
    \lambda=\frac{49}{9}-8=\frac{49-72}{9}=-\frac{23}{9}.
\]
% FIX p.65: the notes write this reduction as $(49-8)/7=-23/9$.
% The sum of roots of \eqref{eq:pq-cubic} is $49/9$, so $\lambda=49/9-8=-23/9$.
% The value $-23/9$ written underneath is the correct root; $(49-8)/7=41/7$ is not.
The line at this abscissa has
\begin{align*}
    y
    &=\frac{7}{3}\left(-\frac{23}{9}\right)-\frac{1}{3}
     =\frac{-161-9}{27}
     =-\frac{170}{27}.
\end{align*}
% FIX p.65: the notes record the intersection ordinate as $-170/9$ and the sum as $(-23/9,\,170/9)$.
% The line is $y=(7x-1)/3$, so $y=(-170/9)/3=-170/27$.
Reflecting across the $x$-axis,
\[
    P+Q=\left(-\frac{23}{9},\,\frac{170}{27}\right).
\]
\end{example}

\lectureheading{Lecture 22 (03/11/26)}

To double a point, replace the chord by the tangent. Let the tangent line to $E$ at $P$ meet $E$ again at $R$. Then $P+P$ is the reflection of $R$ across the $x$-axis.

\begin{example}[Doubling]
Continue with $y^2=x^3-15x+18$ and $P=(7,16)$. Differentiating implicitly,
\[
    2yy'=3x^2-15,
    \qquad
    y'=\frac{3x^2-15}{2y}.
\]
The slope of the tangent at $P$ is
\[
    \frac{3\cdot 7^2-15}{2\cdot 16}
    =\frac{3\cdot 49-15}{32}
    =\frac{132}{32}
    =\frac{33}{8}.
\]
The tangent line is
\[
    y=\frac{33}{8}(x-7)+16
    =\frac{33}{8}x-\frac{103}{8}.
\]
% --- page 66 ---
For the third intersection, substitute into the curve:
\[
    \left(\frac{33x-103}{8}\right)^2=x^3-15x+18.
\]
Expanding $(33x-103)^2=1089x^2-6798x+10609$ and $64(x^3-15x+18)=64x^3-960x+1152$, rearranging yields
\[
    64x^3-1089x^2+5838x-9457=0,
\]
or
\[
    x^3-\frac{1089}{64}x^2+\frac{2919}{32}x-\frac{9457}{64}=0.
\]
The point $P$ is a point of tangency, so $x=7$ is a double root. The third root therefore satisfies
\[
    7+7+x=\frac{1089}{64},
    \qquad
    x=\frac{1089}{64}-14
     =\frac{1089-896}{64}
     =\frac{193}{64}.
\]
% FIX p.66: the third root is written as $(1089-14)/64=193/64$.
% That quotient is $1075/64$. The sum of roots gives $1089/64-14=193/64$, which is the value used below.
The tangent at this abscissa has
\begin{align*}
    y
    &=\frac{33}{8}\cdot\frac{193}{64}-\frac{103}{8}
     =\frac{33\cdot 193-103\cdot 64}{512}
     =\frac{6369-6592}{512}
     =-\frac{223}{512}.
\end{align*}
Reflecting across the $x$-axis,
\[
    2P=P+P=\left(\frac{193}{64},\,\frac{223}{512}\right).
\]
\end{example}

\paragraph{Inverses and the identity.}
The same rule has to add a point to its reflection $(x,-y)$, and a point to the point at infinity.

The line through $(x,y)$ and $(x,-y)$ is vertical, so it meets $E$ at no further affine point. Declare that the point at infinity lies on every vertical line, and write $\mathcal{O}$ for that point. The vertical line then meets $E$ at $(x,y)$, $(x,-y)$, and $\mathcal{O}$, and the group law gives
\[
    (x,y)+(x,-y)=\mathcal{O}.
\]
In the other direction, the vertical line through $P=(x,y)$ meets $E$ again at $P'=(x,-y)$. Reflecting $P'$ across the $x$-axis returns $P$, so
\[
    P+\mathcal{O}=P.
\]

% --- page 67 ---

Thus $P+P'=\mathcal{O}$ with $P'=(x,-y)$. The point $(x,-y)$ is the inverse of $P$, written $-P$.

\begin{keybox}
\begin{theorem}[Group law]
\label{thm:group-law}
Let $E$ be an elliptic curve as above. The addition just defined has the following properties.
\begin{enumerate}[leftmargin=1.6em,label=\textup{(\arabic*)}]
    \item $P+\mathcal{O}=\mathcal{O}+P=P$.
    \item $P+(-P)=\mathcal{O}$.
    \item $P+Q=Q+P$.
    \item $(P+Q)+R=P+(Q+R)$.
\end{enumerate}
Associativity needs a separate rigorous proof; it is not checked from the pictures.
\end{theorem}
\end{keybox}

\subsection{The addition algorithm}

Let $E\colon y^2=x^3+Ax+B$, and let $P_1,P_2$ be points of $E$.
\begin{enumerate}[leftmargin=1.6em,label=\textup{(\arabic*)}]
    \item If $P_1=\mathcal{O}$, then $P_1+P_2=P_2$. If $P_2=\mathcal{O}$, then $P_1+P_2=P_1$.
    \item Otherwise write $P_1=(x_1,y_1)$ and $P_2=(x_2,y_2)$. If $x_1=x_2$ and $y_1=-y_2$, then $P_1+P_2=\mathcal{O}$.
    \item Otherwise define $\lambda$, $x_3$, and $y_3$ as follows, and set $P_1+P_2=(x_3,y_3)$.
\end{enumerate}

% --- page 68 ---

In case~(3),
\begin{equation}
    \lambda=
    \begin{cases}
        \dfrac{y_2-y_1}{x_2-x_1} & \text{if }P_1\neq P_2,\\[1.1em]
        \dfrac{3x_1^2+A}{2y_1} & \text{if }P_1=P_2,
    \end{cases}
    \label{eq:lambda}
\end{equation}
and
\begin{align}
    x_3&=\lambda^2-x_1-x_2, \label{eq:x3}\\
    y_3&=\lambda(x_1-x_3)-y_1. \label{eq:y3}
\end{align}
% FIX p.68: the formula line writes $y_3=\lambda(x_1-x_3)+y_1$.
% The proof just below derives the minus sign, which is the group law.
Then $P_1+P_2=(x_3,y_3)$.

\begin{proof}
In case~(3), $\lambda$ is the slope of the line through $P_1$ and $P_2$ (the tangent line, when the points coincide). That line is
\[
    y=\lambda x+\nu,
    \qquad
    \nu=y_1-\lambda x_1.
\]
Substitute into the curve:
\begin{align*}
    (\lambda x+\nu)^2&=x^3+Ax+B,\\
    x^3-\lambda^2 x^2+(A-2\lambda\nu)x+(B-\nu^2)&=0.
\end{align*}
The three roots are the abscissae of the intersections, so
\[
    \lambda^2=x_1+x_2+x_3,
    \qquad
    x_3=\lambda^2-x_1-x_2.
\]
The third intersection has ordinate $\lambda x_3+\nu$. Reflecting it across the $x$-axis gives the sum:
\begin{align*}
    y_3
    &=-(\lambda x_3+\nu)
     =-\lambda x_3-\nu
     =-\lambda x_3-(y_1-\lambda x_1)
     =\lambda(x_1-x_3)-y_1.
\end{align*}
\end{proof}

With this addition, the points of an elliptic curve form a group. The identity is the point at infinity $\mathcal{O}$.

\paragraph{Integer points.}
How many points with integer coordinates does $E$ have? For $A,B\in\Z$, Siegel's theorem says there are only finitely many integer points. Baker's method gives a bound on the number of those points.

% --- page 69 ---

\subsection{Rational points}

The rational points are much more structured than a finite set.

\begin{theorem}[Mordell--Weil]
The group $E(\mathbb{Q})$ of rational points, including $\mathcal{O}$, is a finitely generated abelian group:
\[
    E(\mathbb{Q})\;\cong\; G\times \Z^{r}.
\]
Here $G$ is a finite abelian group and the integer $r\geq 0$ is the \emph{rank} of the elliptic curve.
\end{theorem}

The rank is the quantity one most wants to understand for an elliptic curve over $\mathbb{Q}$. Bhargava's theorem says that a positive proportion of elliptic curves over $\mathbb{Q}$ have rank $0$, and a positive proportion have rank $1$.

If $r=0$, then $E(\mathbb{Q})$ itself is a finite abelian group. Mazur's theorem lists $15$ explicit possibilities for $G$. Those groups are too small to use for cryptography.

The Lutz--Nagell theorem says that, on the model $y^2=x^3+Ax+B$ with $A,B\in\Z$, every torsion point is an integer point.
% FIX p.69: the notes write ``Lutz-Nagelli''; the result is the Lutz--Nagell theorem.

\subsection{Elliptic curves over finite fields}

Now take an odd prime $p$, and an equation
\[
    y^2=x^3+Ax+B
    \qquad\text{with}\qquad
    4A^3+27B^2\not\equiv 0\pmod{p}.
\]
% --- page 70 ---
The set of solutions in $\Fp\times\Fp$, together with the point $\mathcal{O}$ at infinity, is a group under the same addition law. For each fixed $x$ one has to solve $y^2=x^3+Ax+B$ in $\Fp$.

\begin{example}[A curve over $\mathbb{F}_{13}$]
\label{ex:f13}
Let $E\colon y^2=x^3+3x+8$ over $\mathbb{F}_{13}$. The group is
\begin{align*}
    E(\mathbb{F}_{13})
    =\{&
        \mathcal{O},\;
        (1,5),\;(1,8),\;
        (2,3),\;(2,10),\\
        &(9,6),\;(9,7),\;
        (12,2),\;(12,11)
    \}.
\end{align*}
The same chord-and-tangent formulae work in $\Fp$. The notes ask whether these points carry a group structure; the answer is that one mimics the same addition law. Take $P=(9,7)$ and $Q=(1,8)$. Then
\[
    \lambda
    =\frac{8-7}{1-9}
    =\frac{1}{-8}.
\]
Modulo $13$ one has $-8\equiv 5$, and $5\cdot 8=40\equiv 1$, so $5^{-1}\equiv 8$ and $\lambda\equiv 8$. Then
\begin{align*}
    x_3
    &=\lambda^2-x_1-x_2
     =8^2-9-1
     =64-10
     =54
     \equiv 2\pmod{13},
\end{align*}
since $54-4\cdot 13=2$, and
\begin{align*}
    y_3
    &=\lambda(x_1-x_3)-y_1
     =8(9-2)-7
     =8\cdot 7-7
     =49
     \equiv 10\pmod{13},
\end{align*}
since $49-3\cdot 13=10$. (Equivalently $8\cdot 7-7=7\cdot 7=49$.) Therefore
\[
    P+Q=(2,10),
\]
which is one of the points listed above.
% FIX p.70: a sign in the handwritten line $8^2-9-1=64+10=54$ does not match $64-10=54$. The value $54\equiv 2$ used for $x_3$ is the correct one.
\end{example}

% --- page 71 ---

Doubling the same point $P=(9,7)$,
\begin{align*}
    \lambda
    &=\frac{3x_1^2+A}{2y_1}
     =\frac{3\cdot 9^2+3}{2\cdot 7}
     =\frac{246}{14}.
\end{align*}
Here $246=19\cdot 13-1\equiv -1\pmod{13}$ and $14\equiv 1\pmod{13}$, so $\lambda\equiv -1$. Then
\begin{align*}
    x_3
    &=\lambda^2-x_1-x_2
     =1-9-9
     =-17
     \equiv 9\pmod{13},\\
    y_3
    &=\lambda(x_1-x_3)-y_1
     =-1(9-9)-7
     =-7
     \equiv 6\pmod{13}.
\end{align*}
Hence $2P=(9,6)$.

\paragraph{The size of $E(\Fp)$.}
There are $p$ choices for $x$, and each gives at most two choices for $y$. Counting $\mathcal{O}$ as well,
\[
    \bigl|E(\Fp)\bigr|\leq 2p+1.
\]
About half of the values of $x$ make $x^3+Ax+B$ a quadratic residue, and about half make it a quadratic non-residue. The residue case contributes two points and the non-residue case contributes none, so one expects about $p$ affine points. Together with $\mathcal{O}$,
\[
    \bigl|E(\Fp)\bigr|\;\approx\; p+1.
\]
% FIX p.71: the notes equate $(p\cdot 2+1)/2$ with $p+1$.
% That quotient is $p+1/2$. The count intended by the residue heuristic is $(p/2)\cdot 2+1=p+1$.

\begin{keybox}
\begin{theorem}[Hasse bound]
\[
    \bigl|E(\Fp)\bigr|=(p+1)-t_p
    \qquad\text{with}\qquad
    |t_p|\leq 2\sqrt{p}.
\]
The bound is tight.
\end{theorem}
\end{keybox}
% FIX p.71: the notes spell the name ``Haase''; the bound is Hasse's.

% --- page 72 ---

The integer $t_p=(p+1)-\bigl|E(\Fp)\bigr|$ is called the \emph{trace of Frobenius}.
% FIX p.72: the notes write $t_p=\bigl|E(\Fp)\bigr|-(p+1)$, the negative of the trace appearing in the Hasse formula $|E(\Fp)|=(p+1)-t_p$ on the previous page.

\subsection{The elliptic-curve discrete logarithm}

\begin{keybox}
\begin{definition}[EC-DLP]
Let $E$ be an elliptic curve over $\Fp$, and let $P,Q\in E(\Fp)$. The elliptic-curve discrete logarithm problem asks for an integer $n$ such that
\[
    nP=Q,
\]
that is,
\[
    \underbrace{P+P+\cdots+P}_{n\text{ times}}=Q.
\]
The integer $n$ is called the discrete logarithm of $Q$ with respect to $P$, written $\log_P(Q)$. This is the additive form of the index $\ind$ of Lecture~12.
\end{definition}
\end{keybox}

\begin{remark}
As for the general discrete logarithm in Lecture~12, the equation $nP=Q$ can be solved if and only if $Q$ lies in the cyclic subgroup $\langle P\rangle$.
% FIX p.72: the notes spell ``solved'' as ``sowed''.
The solution is not unique: it is determined only modulo the order of $P$.
\end{remark}

\begin{example}[A curve over $\mathbb{F}_{73}$]
Let $y^2=x^3+8x+7$ over $\mathbb{F}_{73}$, and set
\begin{align*}
    P&=(32,53),\\
    Q&=(39,17)=11P,\\
    R&=(35,47)=37P.
\end{align*}
The order of $P$ is $41$. Since $11+37=48\equiv 7\pmod{41}$,
\[
    Q+R=11P+37P=48P=7P=(37,41).
\]
\end{example}

Alice, in the protocols built on this group, needs to compute $nP$. Adding $P$ to itself one step at a time produces the sequence $P,\,2P,\,3P,\,4P,\,\ldots$, which takes $n$ additions.

% --- page 73 ---

\paragraph{Double-and-add.}
Write the binary expansion
\[
    n=n_0+n_1\cdot 2+\cdots+n_r\cdot 2^r,
    \qquad
    n_i\in\{0,1\},
    \quad n_r=1.
\]
Compute the successive doubles
\[
    Q_0=P,
    \quad
    Q_1=2Q_0,
    \quad
    Q_2=2Q_1,
    \quad
    \ldots,
    \quad
    Q_r=2Q_{r-1}.
\]
Then $Q_i=2^i P$, and
\[
    nP=\sum_{i=0}^{r} n_i Q_i.
\]
Building the $Q_i$ costs $r$ doublings, and summing the terms with $n_i=1$ costs at most $r$ further additions, so $nP$ is obtained in at most $2r$ group operations. Because $n_r=1$,
\[
    2^r\leq n<2^{r+1},
    \qquad
    r\leq \log_2 n<r+1.
\]
% FIX p.73: the notes write the stricter inequality $r<\log_2 n<r+1$. Equality $r=\log_2 n$ holds whenever $n$ is a power of $2$.

\clearpage
\phantomsection
\addcontentsline{toc}{part}{Summary}
\section*{Summary of Lectures 13--22}
\vspace{0.4em}

\[
\boxed{
\begin{array}{c}
\text{Shanks / Pohlig--Hellman}
\longrightarrow
\text{RSA}
\longrightarrow
\text{Miller--Rabin, AKS}
\\[4pt]
\longrightarrow
\text{Pollard $\rho$, Pollard $p-1$}
\longrightarrow
\text{signatures, quadratic sieve}
\\[4pt]
\longrightarrow
\text{Goldwasser--Micali, index calculus}
\longrightarrow
\text{elliptic curves / EC-DLP}.
\end{array}
}
\]

Main facts:
\[
    g^{x}=h
    \text{ in }
    O\bigl(\sqrt{N}\log N\bigr)
    \text{ steps (Shanks)},
    \qquad
    \sum_{i}e_{i}\,S(p_{i})+O(\log N)
    \text{ (Pohlig--Hellman)},
\]
\[
    N=pq,
    \qquad
    de\equiv 1\pmod{\varphi(N)},
    \qquad
    m\equiv C^{d}\pmod{N},
\]
\[
    a^{q}\equiv 1
    \text{ or }
    a^{2^{j}q}\equiv -1\pmod{p},
    \qquad
    (x+a)^{n}\equiv x^{n}+a
    \quad\text{in }(\Z/n\Z)[x],
\]
\[
    \gcdop(a^{L}-1,N)=p
    \text{ when }(p-1)\mid L\text{ and }a^{L}\not\equiv 1\pmod{q},
\]
\[
    a^{2}\equiv b^{2}\pmod{N}
    \text{ gives }
    \gcdop(a\pm b,N),
\]
\[
    s_{1}\equiv g^{k}\pmod{p},
    \qquad
    s_{2}\equiv (D-as_{1})k^{-1}\pmod{p-1},
    \qquad
    A^{s_{1}}s_{1}^{s_{2}}\equiv g^{D}\pmod{p},
\]
\[
    L(x)=\exp\bigl(\sqrt{\log x\cdot\log\log x}\bigr),
    \qquad
    \psi(x,B)=x\,u^{-u(1+o(1))},
    \qquad
    u=\frac{\log x}{\log B},
\]
\[
    \Bigl(\frac{c}{p}\Bigr)=(-1)^{m},
    \qquad
    \ind_{g}(h)
    \equiv
    -k+\sum_{\ell}e_{\ell}\,\ind_{g}(\ell)
    \pmod{p-1},
\]
\[
    x_{3}=\lambda^{2}-x_{1}-x_{2},
    \qquad
    y_{3}=\lambda(x_{1}-x_{3})-y_{1},
    \qquad
    \bigl|E(\Fp)\bigr|=(p+1)-t_{p},
    \quad
    |t_{p}|\le 2\sqrt{p},
\]
and
\[
    nP=Q
    \quad\text{in }O(\log n)\text{ additions}.
\]
